\clearpage
\newpage

\section{Prompt}\label{appendix:prompt}

In this section, we provide the prompt used for invoking summary tools for context summarization within the ReSum paradigm. We intentionally omit explicit instructions asking the summary tool to list current information gaps or provide clear action plans. This design choice aims to avoid two potential issues: (1) distracting the summary tool from its primary task of consolidating key information, and (2) trapping agents in cycles of repeated self-verification due to forced specification of gaps. Remarkably, even without such explicit constraints, we observed that the summary tool retains \textbf{the emergent capability to intuitively identify information gaps and suggest next-step plans when necessary}, thereby balancing summarization fidelity with strategic reasoning.

\begin{tcolorbox}[colback=gray!5, colframe=black, boxrule=1pt, arc=2pt, left=5pt, right=5pt]

\textbf{\textcolor{blue}{Prompt for Context Summarization}} \vspace*{5pt}

You are an expert at analyzing conversation history and extracting relevant information. Your task is to thoroughly evaluate the conversation history and current question to provide a comprehensive summary that will help answer the question.

Task Guidelines: 
\begin{enumerate}
    \item \textbf{Information Analysis}
      \begin{itemize}
          \item Carefully analyze the conversation history to identify truly useful information.
          \item Focus on information that directly contributes to answering the question.
          \item Do NOT make assumptions, guesses, or inferences beyond what is explicitly stated in the conversation.
          \item If information is missing or unclear, do NOT include it in your summary.
      \end{itemize}
    \item \textbf{Summary Requirements}
      \begin{itemize}
          \item Extract only the most relevant information that is explicitly present in the conversation.
          \item Synthesize information from multiple exchanges when relevant. Only include information that is certain and clearly stated in the conversation.
          \item Do NOT output or mention any information that is uncertain, insufficient, or cannot be confirmed from the conversation.
          
      \end{itemize}
   \item \textbf{Output Format} Your response should be structured as follows: 

     \texttt{\textbf{\textcolor{blue}{<summary>}}} 

     \begin{itemize}
         \item Essential Information: [Organize the relevant and certain information from the conversation history that helps address the question.]
         % \item Information Gaps: [List the specific missing pieces that prevent a complete answer.]
         % \item Next-Step Plan: [Propose 1–3 precise actions (e.g., search queries, URLs to visit, sections to read) to resolve the gaps.]
     \end{itemize}

     \texttt{\textbf{\textcolor{blue}{</summary>}}} 
\end{enumerate}

Strictly avoid fabricating, inferring, or exaggerating any information not present in the conversation. Only output information that is certain and explicitly stated.

Question \texttt{ \{Question\} }

Conversation \texttt{ \{Conversation History\} }

Please generate a comprehensive and useful summary.
\end{tcolorbox}


After summary generation, we concatenate the initial question and the summary as a new formatted query for the agent to continue reasoning.

\begin{tcolorbox}[colback=gray!5, colframe=black, boxrule=1pt, arc=2pt, left=5pt, right=5pt]
\textbf{\textcolor{blue}{Prompt for Summary-conditioned Reasoning}} \vspace{5pt}

Question \texttt{ \{Question\} }

Below is a summary of the previous conversation. This summary condenses key information from earlier steps, so please consider it carefully. Assess whether the summary provides enough information to answer the question and use it as the basis for further reasoning and information gathering to answer the question.

Summary: \texttt{\{Summary\}}
\end{tcolorbox}

